% !PRINTED 143
\weberpage{143}
\begin{center}
\S.\,36.

\gesperrt{Limits of the index of a divisor of the symmetric
permutation group.}
\end{center}

We conclude these considerations with the proof of a
theorem on permutation groups which, through the difficulty
that its proof at first presented, has attained a certain celebrity,
and which is of importance for the judgement of algebraic
questions\,$^{1}$).

It concerns the symmetric permutation group
 $P$ of $n$ figures and its divisors of smallest
possible index. We know that the group $P$ always has a
divisor of index 2, namely the alternating group.
Besides this, a divisor of index $n$ is also known, which comprises all
permutations of $P$ that leave one figure unchanged,
and which is therefore intransitive.

First of all the following theorem holds:

I. \gesperrt{The index of an imprimitive divisor of $P$ is
always greater than $n$, and the index of an intransitive
 divisor is equal to or greater
than $n$, and equal to $n$ only when the divisor
leaves one figure at rest and permutes the remaining
$n - 1$ figures in all possible ways.}

Let us assume that $Q$ is an imprimitive divisor of $P$
of index $j$, and that there are $r$ systems of imprimitivity of
$s$ figures each, so that $n = rs$. A number which the degree of the
group $Q$ certainly cannot exceed is obtained if we
count all the permutations in each single one of the $r$ systems and then
also all the permutations of the systems. The degree
of $Q$ is therefore less than or equal to the product $[\Pi(s)]^r\,\Pi(r)$,
if for every integer $n$
\[
\Pi(n) = 1.2.3\dots n
\]
%CONT>>

$^{1}$) \gesperrt{Bertrand}, Journal de Mathématiques, Tome XV (1845). \gesperrt{Serret},
Algèbre supér., Section IV, Chapitre III. C. \gesperrt{Jordan}, Traité des substitutions,
p.\,67. \gesperrt{Netto}, Substitutionentheorie, Capitel VI. Crelle's Journ.,
Bd. 100.
